\scanpage{130}
\mspage{1}
% archive sheet number: 129
12.7.86 \hfill CF IX bis

% margin note written slanted, rising along the upper left edge.
\margin{Comparer CF $\overline{\underline{\unsure{\mathrm{VIII}|\mathrm{XIII}}}}$ p. 46 ff. et aussi p. 27 ff.}
\uline{Morphismes d'ateliers}

Je suis intéressé par trois types de situations, que je voudrais pouvoir exprimer par une notion de morphismes

1\textsuperscript{o}) Cas d'un atelier \inserted{$A$} et d'un sous-atelier $A'$, i.e. (cf CF $\overline{\underline{\mathrm{IX}}}$ p. 27) une partie $A'$ satisfaisant la condition
\[
\begin{aligned}
&\text{\uline{Sous at}}\quad \text{Si }X\overset{\circ}{\ll}Y'\triangleleft Y,\quad\text{avec }X,Y\in A',\ \text{alors}\\
&Y'\in A'.\quad\unsure{!}
\end{aligned}
\]

Cela implique que, muni des relations induites par $A$ $\unlhd$, $\overset{\circ}{\ll}$, $|\unsure{\circ|\bullet}|$, $A'$ est un atelier, et que de plus, la relation $\ll_{A'}$ est induite par $\ll_A$ (ce qui est \unsure{assuré} précisément par Sous at). Je rappelle les exemples principaux

\uline{Ex 1} Soit \unsure{\textsuperscript{*}} $\unsure{\underline{\varphi}|\underline{\Phi}}\inserted{\subset A}$ une préfigure de $A$, i.e. une partie telle que $\overline{\underline{\varphi}}^{(\unlhd)}$ soit une figure, i.e. telle que $X,\deleted{\unsure{X}}\unsure{Y}\in\underline{\varphi}\Rightarrow X\unsure{\asymp|\bowtie}Y$. Alors $\operatorname{Omb}^{\circ}(\underline{\varphi})=\{X\in A\mid X\overset{\circ}{\ll}\underline{\varphi}\}$ est un sous-atelier. (loc. cit)

\uline{Ex} \uline{1 bis} $\underline{\varphi}$ une \deleted{\unsure{pré}}figure de $A$, alors $\underline{\varphi}$ lui-même\inserted{!} est un \unsure{sous-atelier}, (qui est une \uline{maquette} pour la structure induite). (En effet, si $X\overset{\circ}{\ll}Y'\unlhd Y$, avec $X,Y\in\underline{\varphi}\inserted{\subset\overline{\underline{\varphi}}}$, on a $Y'\in\overline{\underline{\varphi}}$, \unsure{mais aussi} $X=Y'$, d'où $Y'\in\underline{\varphi}$.)

